\documentclass[12pt]{article}
\usepackage{amsmath}
\usepackage{amssymb}  % Provides \mathbb and other symbols
\usepackage{geometry}
\usepackage{graphicx}
\usepackage{subcaption}
\usepackage{enumitem}
\usepackage{listings}
\usepackage{hyperref}
\usepackage{ulem}
\usepackage{xcolor}
\usepackage{amsthm}
\usepackage{tcolorbox}  % For colored boxes
\usepackage[table]{xcolor} % For coloring table cells
\usepackage{array}         % For better table alignment
\usepackage{mdframed}
\usepackage{cancel}

% Set the top margin
\geometry{top=1in} % You can adjust the measurement as needed


\begin{document}

\begin{titlepage}
\centering
\vspace*{\fill} % Adds vertical space to center the title block vertically

{\Huge Algebra} \\[0.5cm] % Title 2
{\Huge Home Work 3}\\[1.5cm] % Title 3
{\Large Student: Mark Roberts}\\[0.5cm] % Student ID
{\Large \today} % Date, \today adds today's date

\vspace*{\fill}
\end{titlepage}
\newpage
\noindent
\section*{\textbf{Question 3.1.1}}
\textbf{Part 1}: \\
Let $\phi : G \to H$ be a homomorphism between the two groups $G$ and $H$.  Let $E \le H$ be a subgroup of H.  Prove that: $\phi^{-1}(E) \le G$ (i.e. $\phi^{-1}(E)$ is a subgroup of $G$).\\
\\
Let $K=\phi^{-1}(E) \subset G$.  We just need to show that $K$ is a subgroup of $G$.  We apply the usual subgroup tests:\\
\begin{itemize}[noitemsep, topsep=0pt]
    \item \textbf{Check $K$ not empty}:\\
    From the properties of a homomorphism we know that $1_G \in K$ (and hence $K$ not empty).\\
    \item \textbf{Check $K$ contains inverses}:\\
    $k \in K \implies \phi(k) \in E \implies \phi(k)^{-1} \in E \implies \phi(k^{-1}) \in E \implies k^{-1} \in K$.\\
    These steps follow as $E$ itself is a group and $\phi$ a homomorphism. \\
    \item \textbf{Check $K$ stable}:\\
    Let $a, b \in K \implies \phi(a), \phi(b) \in E \implies \phi(a)\phi(b) = \phi(ab) \in E \implies ab \in K$.
\end{itemize}
\qed \\
\\
\textbf{Part 2}: \\
If $E \trianglelefteq H$ prove that $\phi^{-1}(E) \trianglelefteq G$.  Deduce that: $Ker(\phi) \trianglelefteq G$. \\
\\
From part (1) we know that $E \le H \implies \phi^{-1}(E) \le G$; so we just need to confirm that $E$ normal in $H$ implies that $K = \phi^{-1}(E)$ normal in $G$.  Now take any $g \in G$ and $k \in K$.  Then:
\begin{align*}
    &\phi(gkg^{-1}) = \phi(g)\phi(k)\phi(g^{-1}) = \phi(g)\phi(k)\phi(g)^{-1} \;\;\;\;\;\;\;\;\;\; \text{(since $\phi$ is a homomorphism)} \\
    \implies &\phi(gkg^{-1}) \in E  \;\;\;\;\;\;\;\;\;\;\;\;\;\;\;\;\;\;\;\;\;\;\;\;\;\;\;\;\;\;\;\;\;\;\;\;\;\;\;\;\;\;\;\;\;\;\;\;\;\;\;\;\;\;\;\;\;\;\; \text{(since $\phi(k) \in E$ and $E$ is normal)} \\
    \implies &gkg^{-1} \in \phi^{-1}(E) = K
\end{align*}
Hence, $K=\phi^{-1}(E)$ must be normal in $G$.

\newpage
\section*{\textbf{Question 3.1.3}}
Given that $A$ is an abelian group and that $B \le A$.  Prove that $A/B$ is an abelian group. \\
\\
Since $A$ is abelian we know that $B$ is a normal subgroup and that $A/B$ is a group. To show that it is itself abelian let us first look at the definition of the group $A/B$:
\begin{align*}
    A/B &= \{aB \; | \; a \in A \}
\end{align*}
We now just need to show that the binary operation defined on $A/B$ is commutative.\\
\\
Let $\alpha B$ and $\beta B$ (where $\alpha, \beta \in A$) be two elements of this set.  Then:
\begin{align*}
    (\alpha B)(\beta B) &= \{(\alpha b_1)(\beta b_2) \; | \; b_1, b_2 \in B \} \\
        &= \{\alpha b_1 \beta b_2 \; | \; b_1, b_2 \in B \} \\
        &= \{\beta b_1 \alpha b_2 \; | \; b_1, b_2 \in B \} \;\;\;\;\;\;\;\;\;\;\;\;\;\;\;\; \text{(since $A$ is abelian)} \\
        &= \{(\beta b_1)(\alpha b_2) \; | \; b_1, b_2 \in B \} \\
        &= (\beta B)(\alpha B)
\end{align*}
\qed

\newpage
\section*{\textbf{Question 3.1.5}}
Let $G$ be a group and $G/N$ a quotient group (and hence $N$ is normal).  Show that the order of $gN \in G/N$ is the smallest power of $g$ such that $g \in N$.\\
\\
Firstly, the order of $gN \in G/N$ is the smallest $n \in \mathbf{N}$ such that: $(gN)^n = N$.  Now, this is equivalent to: $g^nN = N$.  But this last equation only holds when $g^n \in N$.  So, clearly, we have that: $|gN|$ is the smallest positive integer such that $g^n \in N$.\\
\qed

\newpage
\section*{\textbf{Question 3.1.9}}
Define: $\phi: \mathbf{C}^{\times} \to \mathbf{R}^{\times}$; by $\phi(a+bi) = a^2 + b^2$.\\
\\
\textbf{1. Proof $\phi$ is a homomorphism}:\\
We have two properties that we must prove:
\begin{itemize}[noitemsep, topsep=0pt]
    \item Check $\phi$ maps identity to identity.\\
    Now $1_{\mathbf{C}^{\times}}=1$ and $1_{\mathbf{R}^{\times}}=1$.  And we see that: $\phi(1)=1$ - so this property holds.\\
    \item Check $\phi(ab) = \phi(a)\phi(b)$:\\
    Let $a=a_1 + a_2 i$ and $b=b_1 + b_2 i$.  Then:
    \begin{align*}
      \phi(ab) &= \phi((a_1 + a_2 i)(b_1 + b_2 i)) = \phi(a_1 b_1 - a_2 b_2 + (a_1 b_2 + a_2 b_1)i) \\
      &= [a_1 b_1 - a_2 b_2]^2 + [a_1 b_2 + a_2 b_1]^2 \\
      &= [(a_1b_1)^2 + (a_2b_2)^2 - 2(a_1b_1a_2b_2)] + [(a_1b_2)^2 + (a_2b_1)^2 + 2(a_1b_2a_2b_1)] \\
      &= (a_1^2 + a_2^2)(b_1^2 + b_2^2) \\
      &= \phi(a_1 + a_2 i) \phi(b_1 + b_2 i) \\
      &= \phi(a) \phi(b)
    \end{align*}
\end{itemize}
These two properties holding demonstrates that $\phi$ is indeed a homomorphism.\\
\\
\\
The kernel of this homomorphism are all those complex numbers in $\mathbf{C}^{\times}$ whose modulus is 1.
This is simply those complex numbers of modulus 1 - i.e. the circle of complex numbers of radius 1 around the origin.

\newpage
\section*{\textbf{Question 3.1.10}}


\newpage
\section*{\textbf{Question 3.1.14}}

\newpage
\section*{\textbf{Question 3.1.20}}


\newpage
\section*{\textbf{Question 3.1.32}}
$Q_8$ is the quaternion group, where:
\begin{align*}
      Q_8 = \{\pm1,\pm i,\pm j,\pm k \} \tag{3.1.32.1} \\
      i^2=j^2=k^2=ijk=-1 \tag{3.1.32.2} \\
      ij=k,\quad jk=i,\quad ki=j  \tag{3.1.32.3}
\end{align*}
1) Prove every subgroup of $Q_8$ is normal. \\
2) For each subgroup find the isomorphism type of its  corresponding quotient.


\end{document}